% Inserted into the Preliminaries / ANN subsection of the Overleaf paper,
% immediately after the LSTM paragraph (Joy, 4 Nov 2026). Red = \JOY.
% The Preliminaries intro sentence was also extended minimally:
%   "... neural network-based forecasting {\JOY and approximation} models,
%    {\JOY including the two feed-forward benchmarks used in the
%    approximation part of the study,} ..."
{\JOY
\paragraph{Multilayer Perceptrons (MLP) and Fourier features.}
A multilayer perceptron with two hidden layers of width \(w\) maps a
scalar input \(t\) to
\begin{equation}
f(t)=W_3\,\tanh\!\bigl(W_2\,\tanh(W_1 t+b_1)+b_2\bigr)+b_3,
\label{eq:mlp}
\end{equation}
with fixed activation functions on the nodes and learnable weights on
the edges, the complement of the KAN arrangement. Gradient-trained MLPs
fit low-frequency components of a target first and high-frequency
detail only slowly, a property known as spectral
bias~\cite{rahaman2019spectral}. A standard remedy is to replace the
input by a random Fourier feature map~\cite{tancik2020fourier}
\begin{equation}
\gamma(t)=\bigl[\sin(2\pi b_1 t),\ldots,\sin(2\pi b_m t),
\cos(2\pi b_1 t),\ldots,\cos(2\pi b_m t)\bigr],\qquad
b_j\sim\mathcal{N}(0,\sigma^2),
\label{eq:fourier}
\end{equation}
drawn once and kept fixed, so that the network operates on
\(\gamma(t)\) rather than on \(t\). Both forms are used as feed-forward
benchmarks for the KAN in the approximation part of the study, following
the practice of comparing KANs and MLPs at matched parameter
counts~\cite{liu2024kan,yu2024kanmlp}.
}
